\documentclass[12pt]{article}
\usepackage{makeidx}
\usepackage[margin=1in]{geometry}  % set the margins to 1in on all sides
\usepackage{graphicx}              % to include figures
\usepackage{amsmath}               % great math stuff
\usepackage{amsfonts}              % for blackboard bold, etc
\usepackage{amsthm}                % better theorem environments
\usepackage{makeidx}               % index
\usepackage[utf8]{inputenc}        % now we have tildes!
\usepackage{wrapfig}               % images
\usepackage{listings}              % Unordered lists
\usepackage{supertabular} 		   % for itemizes with titles

% various theorems, numbered by section

\makeindex

\begin{document}

\begin{titlepage}

\newcommand{\HRule}{\rule{\linewidth}{0.5mm}} % Defines a new command for the
                                    % horizontal lines, change thickness here

\center % Center everything on the page

%-------------------------------------------------------------------------------
%	HEADING SECTIONS
%-------------------------------------------------------------------------------

\textsc{\LARGE Universidad Carlos III de Madrid}\\[1.5cm] % University name
\textsc{\Large Aprendizaje Automático}\\[0.5cm] % Course name
\textsc{\large Computer Science Engineering}\\[0.5cm] % Course title

%-------------------------------------------------------------------------------
%	TITLE SECTION
%-------------------------------------------------------------------------------

\HRule \\[0.4cm]
{ \huge \bfseries Tutorial 2: Introducción a Weka}\\[0.4cm] % Document title
\HRule \\[1.5cm]

%-------------------------------------------------------------------------------
%	AUTHOR SECTION
%-------------------------------------------------------------------------------

\emph{Authors:}\\
Daniel \textsc{Medina García}\\ % Your name
Alejandro \textsc{Rodríguez Salamanca}\\[3cm] % Your name

%-------------------------------------------------------------------------------
%	DATE SECTION
%-------------------------------------------------------------------------------

{\large \today}\\[3cm] % Date, change the \today to a set date if you want to be
                       % precise

%-------------------------------------------------------------------------------
%	LOGO SECTION
%-------------------------------------------------------------------------------

% \includegraphics{Logo}\\[1cm] % Include a department/university logo - this
% will require the graphicx package

%-------------------------------------------------------------------------------

\vfill % Fill the rest of the page with whitespace

\end{titlepage}

\tableofcontents

\newpage

\begin{center}
\section{Los ficheros de datos}
\end{center}

\subsection*{\small ¿Cuántos atributos de entrada tiene el fichero de datos?
¿De qué tipo son?}

\begin{tabular}{rl}
  \texttt{badges} & Este archivo contiene un atributo de entrada de tipo String
  que se corres- \\ &ponde con un nombre, y un atributo de salida, la clase a la
  que pertenece.\\
  \texttt{badges\_plain} & El contenido de este archivo es el mismo que el de
  \texttt{badges}, con la diferen- \\ &cia de que los nombres están previamente
  enumerados.
\end{tabular}

\subsection*{\small ¿Podría un algoritmo de aprendizaje automático identificar
esa función con los datos que hay en ese fichero? ¿Por qué?}

No. El fichero contiene tan solo un campo de datos, de tipo enumerado, y cada
posible valor de dicho campo sólo se encuentra en una instancia. Al ser una
clase con dos valores (muchos menos que valores del enumerado) y ser todas las
instancias distintas, toda información que saque un algoritmo de aprendizaje
automático será meramente memorización, sin valor alguno en futuras instancias
con valores distintos a los actuales ya que no existe razonamiento o
computación en la resolución de la clase.

\newpage

\begin{center}
\section{Clasificar con ZeroR}
\end{center}

\subsection*{\small ¿Qué resultado en términos de instancias correctas ofrece el
algoritmo ZeroR?}

ZeroR indica un resultado del 51.0204\% de instancias correctas. Lo único de lo
que nos informa este dato es de que las clases están equilibradas dentro de
nuestro conjunto de datos.

\subsection*{\small ¿Qué ocurre si se selecciona otro algoritmo de clasificación
permitido para ese conjunto de datos?}

Utilizando meta / MultiScheme obtenemos exactamente el mismo resultado que con
ZeroR. Esto es lógico pues no se puede mejorar el resultado al carecer de
información suficiente para razonar sobre los datos.

\subsection*{\small ¿Cuáles son las diferencias al repetir los pasos anteriores
con el otro fichero, \texttt{badges\_plain}?}

Los mismos, pues aunque cambiemos el tipo de campo de enumerado a String se
sigue sin tener información para razonar sobre los datos.

\subsection*{\small ¿Qué ocurre si seleccionamos el algoritmo trees / ID3 en el
segundo fichero?}

Si probamos con ID3 podemos observar que el porcentaje de instancias
correctamente clasificadas es del 100\%. Si usásemos validación cruzada
seguramente los resultados se parecerían más a los obtenidos con ZeroR, ya que
se sigue sin tener información para clasificar más allá de la memorización y
por ello no se tiene un criterio razonado para clasificar nuevas instancias.

\newpage

\begin{center}
\section{Generando nuevos atributos}
\end{center}

\subsection*{\small Propón 6 nuevos atributos y explica por qué los has elegido.}

Al tener como datos de entrada nombres propios, sólo podemos basarnos en la
naturaleza de las palabras que los componen para realizar algún tipo de
razonamiento. Así, proponemos algunos nuevos campos a partir del original que
detallan las características de dichos nombres:

\begin{itemize}
    \item Número de letras del nombre.
    \item Número de vocales.
    \item Número de consonantes.
    \item Número de espacios.
    \item Inicial del nombre o del apellido.
    \item Número de repeticiones de cada una de las letras que aparecen.
\end{itemize}

\subsection*{\small ¿Cuántos atributos tiene el fichero \texttt{badges1} y de
qué tipo son?}

Este archivo contiene atributos que persiguen lo mismo que planteábamos en la
pregunta anterior: descubrir más datos sobre los nombres. Así, el archivo
contiene siete nuevos campos de entrada además del nombre de tipos numérico y
enumerado.

\subsection*{\small ¿Qué otro tipo de información estadística se muestra sobre
los atributos? Tras pulsar el botón ``Visualize all'' indica qué se muestra y
si hay algún atributo que no se visualice.}

En cuanto a la información estadística sobre los atributos, si el atributo es
de tipo nominal, aparece el valor y las veces que está repetido. En cambio, si
es de tipo numérico, Weka nos ofrece el máximo, el mínimo, la media y la
desviación estándar.

``Visualize all'' muestra la correlación entre los valores para cada atributo
y su pertenencia a una u otra clase.


\subsection*{\small Genera un clasificador con ZeroR, ¿qué ocurre? Compara los
resultados con los obtenidos en el ejercicio anterior.}

ZeroR es un clasificador que sólo tiene en cuenta el número de instancias
que pertenecen a cada clase. Es por eso que su éxito al clasificar no varía
respecto al anterior conjunto de datos que no ofrecía información.

\subsection*{\small Genera un clasificador con trees / ID3, ¿qué ocurre? ¿Qué se
podría hacer para solucionar este problema?}

No es posible generar un clasificador con ID3 porque ID3 no acepta datos del
tipo numérico. Así, una posible solución sería convertir el tipo de entradas a
otros tipos de atributo que no sean numerales (discretización).

\newpage

\begin{center}
\section{Clasificar con ID3: resolviendo problemas}
\end{center}

\subsection*{\small ¿Qué información aparece en el desplegable tras abrir la
pestaña \emph{Capabilities} de ZeroR? ¿Qué información proporciona ``more''?}

En el desplegable aparecen los tipos de clase y de atributos que acepta el
algoritmo. Además, indica como información ``Adicional'' el hecho de que ZeroR
no tiene un número mínimo de entradas para su funcionamiento.

``more'' muestra información no funcional sobre el algoritmo. Incluye el
identificador unívoco del algoritmo dentro de Weka, una breve sinopsis
explicando cómo trabaja dicho algoritmo y sus posibles opciones de
configuración.

\paragraph{\small Los atributos de entrada pueden modificarse a través de
tareas de preprocesamiento. En los siguientes pasos vamos a modificar ciertos
atributos de \texttt{badges1} para que pueda clasificarse con ID3.}

\subsection*{\small ¿Qué efecto tiene el filtro de discretización sobre el
conjunto de datos con \emph{bins} igual a 5?}

Discretiza, i.e. separa los valores de un atributo con valores continuos en su
pertenencia a un rango determinado, los atributos previamente de carácter
numérico convirtiéndolos en enumerados con tantos valores distintos como rangos
creados. Creará tantos rangos como le fijemos a través de \emph{bins}.

\subsection*{\small ¿Cuántas instancias clasifica bien ahora ID3 usando un
conjunto de entrenamiento? ¿Qué porcentaje representa? ¿Qué crees que indica
la ``matriz de confusión''? ¿Cuántas instancias de cada tipo se han clasificado
mal?}

ID3 clasifica bien 284 de 294 instancias con \emph{bins} igual a 10, lo que
supone un 96.5986\% de éxito. Tras reducir \emph{bins} a 5, clasifica
correctamente 48 instancias menos (236 de 294), resultando en un 80.2721\% de
precisión. En ambos casos es un resultado considerablemente superior al
obtenido con ZeroR, lo que nos indica que el algoritmo ha conseguido razonar
sobre la información de los atributos para clasificar.

La matriz de confusión desglosa el éxito en la clasificación por cada una de
las clases, y nos deja ver que han sido 9 las instancias erróneamente
clasificadas como pertenecientes a la clase negativa mientras que tan sólo una
instancia fue malinterpretada como positiva.

\subsection*{\small ¿Cuál es la primera instancia del conjunto de entrenamiento
que se clasifica mal cuando se usa la opción de mostrar las predicciones en la
salida? ¿Por qué?}

La séptima, ya que su clase es negativa y nuestro algoritmo la ha clasificado
como positiva con un 57.1\% de probabilidad. La opción marcada hace imprimir una
tabla con tantas filas como instancias a clasificar, y los siguientes campos:

\begin{tabular}{rl}
    \emph{inst\#} & ID de la instancia.\\
    \emph{actual} & Clase a la que realmente pertenece la instancia.\\
    \emph{predicted} & Resultado de la clasificación de nuestro algoritmo.\\
    \emph{error} & Aparecerá un `+' si la instancia fue erróneamente
    clasificada.\\
    \emph{prob. distr.} & Porcentaje de seguridad, para cada clase, de la
    pertenencia de la instan- \\ &cia a dicha clase.
\end{tabular}

\subsection*{\small ¿Cómo se clasificaría la instancia ``Eloisa Figueroa''?
¿Cuáles son los atributos de este nombre? ¿Qué ocurre con los valores de esta
instancia si utilizas el filtro usado anteriormente?}

Siguiendo el árbol, como el nombre tiene una longitud entre 14.2 y 17.8 (15),
menos de 1.4 espacios (1) y menos de 10.8 consonantes (5), la instancia sería
clasificada como negativa.

Otros campos no utilizados son la paridad y la primera letra vocal, de fácil
deducción con valores 0 y 1 respectivamente, la inclusión de puntos en el
nombre (0) y el número de palabras (2).

Si nos fijamos en el reparto de instancias por clase dentro de un mismo valor
para los atributos recién mencionados podemos observar cómo el porcentaje de
cada clase ronda el 50\%, lo que indica que estos campos no añaden información
para el algoritmo, ocurriendo lo mismo con el número de espacios.

Por otro lado, la longitud sí parece jugar un papel importante en la decisión.
Echando un breve vistazo al reparto de clases en los valores discretizados de la
longitud, vemos que según crece la longitud del nombre encontramos más
instancias pertenecientes a la clase negativa, mientras que los nombres más
cortos tienen más probabilidades de pertenecer a la positiva. Parece haber una
cierta correlación con el número de consonantes, pero esta no es tan clara.

\paragraph{\small A continuación modificamos el fichero original introduciendo
el nombre anterior y clasificándolo como positivo, teniendo en cuenta que
si contiene enumerados y se introduce un nuevo valor hay que especificarlo
también en la definición de los valores posibles del enumerado. Después
volvemos a generar el clasificador con ZeroR y training set seleccionado.}

\subsection*{\small ¿Cómo se clasifica la instancia nueva?}

Al haber una ligera mayoría de instancias positivas, ZeroR elige la positiva
como la clase dominante. Así, la nueva instancia es clasificada correctamente
como positiva.

\newpage

\begin{center}
\section{Clasificar con J48 (C4.5)}
\end{center}

\subsection*{\small ¿Cuántas hojas tiene el árbol para \texttt{badges1}
generado con J48 usando conjunto de entrenamiento? ¿Cuántas instancias del
conjunto de entrenamiento clasifica bien? ¿Qué porcentaje representa? ¿Cuántas
instancias de cada tipo se han clasificado mal? ¿Cómo se clasificaría la
instancia ``Eloisa Figueroa''?}

El árbol generado con J48 tiene 20 hojas. Este árbol clasificada correctamente
287 instancias, que equivale a un 97.619\%. Las instancias incorrectamente
clasificadas son 7, 3 en la clase negativa y 4 en la clase positiva.

\subsection*{\small ¿Elegirías este modelo o el generado por ID3? ¿Por qué?}

Respuesta

\subsection*{\small ¿Hemos encontrado la función exacta para generar las
etiquetas? ¿Por qué lo sabes?}

Respuesta

\newpage

\section{Utilizando más atributos con J48 (C4.5)}

\paragraph{\small Es momento de volver a la pestaña de preproceso y generar un
nuevo atributo que calcule el número de vocales. Después se grabará el conjunto
de datos como \texttt{badges1-2}, y con él se construirá un clasificador
con J48. Una vez generado, se anotan el porcentaje de instancias bien
clasificadas y la matriz de confusión, tras lo cual visualizaremos el árbol
generado.}

\subsection*{\small ¿Qué indican los números que aparecen en las hojas del
árbol?}

Respuesta

\subsection*{\small ¿Qué efecto tiene aumentar el valor de ``Jitter'' en la
gráfica que relaciona el nuevo atributo con la clase?}

Respuesta

\subsection*{\small ¿Podrías decir cuál es el rango de vocales más común en el
fichero proporcionado? ¿Se te ocurre algún otro atributo relacionado que pueda
aportar información?}

Respuesta

\subsection*{\small Tras todos estos resultados, ¿qué características o
cualidades crees que deben tener los atributos para maximizar el éxito de los
algoritmos de aprendizaje automático?}

Respuesta

\newpage

\section{Balanceado de datos, selección de características y otros filtros}

\paragraph{\small Para este apartado se carga en Weka \texttt{adult-data}.}

\subsection*{\small ¿Cuántos atributos de entrada tiene este fichero? ¿Cuántas
instancias de entrenamiento?}

Respuesta

\subsection*{\small ¿Qué resultados aparecen ejecutando ZeroR con validación
cruzada? Explica el resultado.}

Respuesta

\subsection*{\small ¿Qué resultados aparecen sólo con las instancias del fichero
\texttt{adult-test.arff}? ¿Son estos resultados comparables a los anteriores?
¿Por qué?}

Respuesta

\subsection*{\small ¿Qué resultados aparecen repitiendo el proceso con J48?
¿Qué porcentaje de mejora ha obtenido respecto a los resultados del ZeroR?}

Respuesta

\subsection*{\small ¿Qué proporción de datos hay de cada clase? ¿Crees que este
porcentaje es apropiado para que un algoritmo de aprendizaje automático aprenda
bien?}

Respuesta

\subsection*{\small ¿Qué ocurre con el atributo de salida cuando se modifican
las instancias de entrenamiento para tener un porcentaje similar entre las dos
clases? ¿Ha descendido el número de ejemplos de entrenamiento?}

Respuesta

\paragraph{\small Tras aplicar este filtro, se evalúan de nuevo con validación
cruzada y supplied test set los algoritmos ZeroR y J48.}

\subsection*{\small ¿Qué resultados dan los algoritmos? ¿Qué resultado crees que
es mejor? ¿Por qué?}

Respuesta

\subsection*{\small ¿Qué atributos consideras inútiles para el algoritmo de
aprendizaje y por qué? Elimina un par de ellos. ¿Qué es lo que ocurre al
repetir la evaluación anterior?}

Respuesta

\subsection*{\small ¿Qué resultados se obtienen aplicando el filtro de
normalización para los atributos numéricos?}

Respuesta

\subsection*{\small Después del procesamiento de datos que has realizado en este
apartado, ¿crees que esto ayuda al proceso de aprendizaje? ¿Por qué? ¿Cuál es
el mejor resultado obtenido? Justifícalo.}

Respuesta

\end{document}
